\section{System Design}
\label{sec:system}

% GOAL: the conceptual design, independent of code. A reader should
% understand HOW it works and WHY each choice was made. ~1.5-2 pages.
% Put an architecture diagram here — it does a lot of work.

\subsection{Design principle}
% The one rule the system is built around: point back to sources,
% don't substitute for thinking. State it, then show how the pipeline
% enforces it.
TODO.

\subsection{Pipeline overview}
% Walk the path: course materials in -> indexing -> student question ->
% retrieval -> filtering to 2-3 references/images -> response that cites
% the exact sources. A figure here is worth a lot.
\begin{figure}[H]
  \centering
  % \includegraphics[width=0.9\linewidth]{figures/architecture.pdf}
  \fbox{\parbox[c][4cm][c]{0.9\linewidth}{\centering TODO: architecture diagram}}
  \caption{System architecture: from course materials to source-pointing responses.}
  \label{fig:arch}
\end{figure}

\subsection{Ingesting course material}
% How slides + textbooks get parsed, chunked, and indexed. Course-scoping.
TODO.

\subsection{Retrieval and filtering}
% How you go from a query to the 2-3 most relevant references/images,
% and how you keep it grounded in what was actually taught in that course.
TODO.

\subsection{The output decision}
% The interesting part (and the thread you carry into your SOPs): deciding
% what to return — hint vs. reference vs. nothing — and where the model's
% judgment can be trusted vs. constrained.
TODO.
